\documentclass[10pt,pdftex,aspectratio=169]{beamer}
\usepackage[utf8]{inputenc}
\usepackage[russian]{babel}
\usepackage{listings}
\usepackage{xcolor}

% Настройка темы презентации (минималистичный строгий стиль)
\usetheme{Madrid}
\usecolortheme{whale}

% Отключение элементов навигации для чистоты слайдов
\setbeamertext{navigation symbols}{}

% Настройки отображения листингов кода
\listingsset{
    basicstyle=\ttfamily\scriptsize,
    breaklines=true,
    frame=single,
    backgroundcolor=\color{gray!5},
    commentstyle=\color{gray},
    keywordstyle=\bfseries\color{blue},
    showstringspaces=false,
    tabsize=2
}

\title[Контроль целостности Linux]{Разработка и внедрение системы контроля целостности файлов и непрерывного мониторинга операционной системы Linux}
\subtitle{Курсовой проект по администрированию ОС Linux}
\author[Поминов, Кудряшов]{Поминов Савелий Сергеевич \and Кудряшов Никита}
\institute[Инженерная группа]{Физический факультет}
\date{\today}

\begin{document}

% Слайд 1: Титульный
\begin{frame}
    \titlepage
\end{frame}

% Слайд 2: Проблематика и целеполагание
\begin{frame}{Проблематика и целеполагание}
    \begin{block}{Актуальность}
        Традиционный анализ системных журналов \textit{постфактум} неэффективен против целевых атак, компрометирующих конфигурационные файлы за доли секунды.
    \end{block}
    \vfill
    \begin{block}{Концепция и Цель}
        Переход к парадигме \textbf{Zero Trust} («никому не доверяй, всегда проверяй»). Разработка и развертывание комплексной микросервисной SIEM-системы реального времени.
    \end{block}
    \vfill
    \begin{itemize}
        \item \textbf{Базовый стек:} Подсистема аудита \texttt{Auditd}
        \item \textbf{Агенты сбора метрик:} \texttt{Auditbeat} (модуль FIM)
        \item \textbf{Центральное аналитическое ядро:} \texttt{ELK Stack} (Elasticsearch, Kibana)
    \end{itemize}
\end{frame}

% Слайд 3: Командное взаимодействие и распределение ролей
\begin{frame}{Распределение зон ответственности}
    \begin{columns}[t]
        \begin{column}{0.5\textwidth}
            \begin{block}{Поминов Савелий Сергеевич}
                \textbf{Tech Lead / Systems Administrator}
                \begin{itemize}
                    \item Проектирование распределенной архитектуры комплекса.
                    \item Низкоуровневая оптимизация и системный харднинг хост-платформы.
                    \item Координация интеграции модулей и финальное тестирование.
                    \item Ведение сквозной технической документации.
                \end{itemize}
            \end{block}
        \end{column}
        \begin{column}{0.5\textwidth}
            \begin{block}{Кудряшов Никита}
                \textbf{DevOps / Security Engineer}
                \begin{itemize}
                    \item Построение контейнерной фабрики на базе Docker Compose.
                    \item Организация изолированных виртуальных сетей.
                    \item Разработка гранулярных правил \texttt{audit.rules}.
                    \item Конфигурирование политик контроля целостности (FIM).
                    \item Автоматизация тестов безопасности.
                \end{itemize}
            \end{block}
        \end{column}
    \end{columns}
\end{frame}

% Слайд 4: Архитектура платформы и механизмы контроля
\begin{frame}{Архитектура платформы и механизмы контроля}
    \begin{itemize}
        \item \textbf{Перехват на уровне ядра:} Фиксация низкоуровневых системных вызовов (\textit{syscalls}): \texttt{execve}, \texttt{openat}, \texttt{chmod}.
        \item \textbf{Мониторинг целостности (FIM):} Непрерывное фоновое сканирование директорий \texttt{/etc/shadow}, \texttt{/etc/passwd}. При изменении — мгновенное вычисление хэш-суммы \textbf{SHA-256} и сравнение с эталоном.
        \item \textbf{Конвейер обработки данных:}
    \end{itemize}
    \begin{center}
        \boxed{\text{Ядро Linux (Auditd)}} $\longrightarrow$ \boxed{\text{Агент Auditbeat (JSON)}} $\longrightarrow$ \boxed{\text{СУБД Elasticsearch}} $\longrightarrow$ \boxed{\text{Панель Kibana}}
    \end{center}
\end{frame}

% Слайд 5: Декларативное развертывание и контейнеризация
\begin{frame}[fragile]{Декларативное развертывание (IaC)}
    \begin{block}{Решение проблемы изоляции агента безопасности в контейнере}
        Использование директивы \texttt{pid: host} и проброс гранулярных привилегий (\texttt{cap\_add}) для прямого взаимодействия с подсистемой аудита хоста.
    \end{block}
\begin{lstlisting}[language=yaml, caption=Фрагмент deploy/docker-compose.yml]
services:
  auditbeat:
    build: context: ./docker/auditbeat
    user: root
    pid: host
    cap_add:
      - AUDIT_CONTROL
      - AUDIT_READ
    volumes:
      - /var/run/docker.sock:/var/run/docker.sock:ro
      - /:/hostfs:ro
\end{lstlisting}
\end{frame}

% Слайд 6: Системный харднинг и оптимизация Linux
\begin{frame}[fragile]{Системный харднинг и оптимизация Linux}
    \begin{block}{Устранение рисков дефицита памяти и отказов СУБД (OOM Killer)}
        Глубокая настройка параметров ядра через конфигурационный файл \texttt{/etc/sysctl.d/99-hardening.conf}.
    \end{block}
\begin{lstlisting}[language=bash, caption=Конфигурация параметров ядра host-системы]
# Расширение областей карты памяти для стабильной работы Elasticsearch
vm.max_map_count=262144

# Запрет создания дампов памяти ядра для защиты чувствительных данных
fs.suid_dumpable=0

# Защита сетевого стека от атак типа TCP SYN Flood
net.ipv4.tcp_syncookies=1
\end{lstlisting}
    \begin{block}{Сетевая изоляция}
        Периметр защищен брандмауэром \texttt{UFW}. Внешний доступ ограничен только веб-интерфейсом Kibana (\texttt{5601}) и защищенным SSH-каналом.
    \end{block}
\end{frame}

% Слайд 7: Испытания системы и симуляция угроз
\begin{frame}[fragile]{Испытания системы и симуляция угроз}
    \begin{block}{Автоматизация контроля работоспособности (test\_security.sh)}
        Сценарий имитирует вектор атаки от имени скомпрометированной учетной записи \texttt{nobody}.
    \end{block}
\begin{lstlisting}[language=bash]
# Попытка несанкционированного чтения базы хэшей паролей
su -s /bin/bash nobody -c "cat /etc/shadow" || true

# Попытка внедрения вредоносной задачи в планировщик cron
su -s /bin/bash nobody -c "touch /etc/cron.d/malicious_job" || true

# Проверка индексации и доставки события в SIEM
curl -s -u elastic:${ELASTIC_PASSWORD} "http://localhost:9200/auditbeat-*/_search?q=user.name:nobody" | grep -q "nobody"
\end{lstlisting}
    \textbf{Результат:} Ядро генерирует прерывание, правила \texttt{audit.rules} маркируют инцидент, и угроза отображается в Kibana менее чем за 1 секунду.
\end{frame}

% Слайд 8: Заключение
\begin{frame}{Заключение}
    \begin{block}{Результаты реализации проекта}
        \begin{itemize}
            \item \textbf{Высокая производительность:} Разработан и протестирован функционирующий прототип легковесной SIEM-системы корпоративного класса.
            \item \textbf{Абсолютная воспроизводимость:} Инфраструктура разворачивается в один клик благодаря строгому следованию методологии IaC.
            \item \textbf{Отказоустойчивость:} Оптимизация параметров ядра Linux позволила полностью исключить аварийные остановки контейнеров СУБД.
            \item \textbf{Автономность:} Платформа успешно прошла стресс-тесты, доказала надежность ведения гранулярного аудита и готова к промышленной эксплуатации.
        \end{itemize}
    \end{block}
\end{frame}

\end{document}